\begingroup
\renewcommand{\arraystretch}{1.45}
\Needspace{44\baselineskip}
\InterpolationArticle{41}

Since all the special cases considered in Articles 29---40 have been completely reduced to the general case in Articles 20---28, all the devices that shorten the calculation in this case, such as those we explained in Articles 25, 26, 27, will also be applicable to them. Hence there will be no need to dwell on this investigation, which the following example, belonging to Articles 39, 40, will bring to a close.

The equation of the center for the new planet \textit{Juno}, using the eccentricity \(0{,}254236\), has been calculated by the indirect method at steps of ten degrees, as follows

\[
\begin{array}{r|r|rrrcr}
\multicolumn{2}{c|}{}&
\multicolumn{5}{c}{\text{Equation of the center}=X}\\
\multicolumn{2}{r|}{\text{Mean anomaly}=x}&-&+&&&\\ \hline
0^{\mathord{0}}&360^{\mathord{0}}&0&&&&0\\
10&350&3^{\mathord{0}}&50'&38''30&\cdot\ \cdot\ \cdot\ \cdot&13838''30\\
20&340&7&38&21{,}47&\cdot\ \cdot\ \cdot\ \cdot&27501{,}47\\
30&330&11&20&8{,}79&\cdot\ \cdot\ \cdot\ \cdot&40808{,}79\\
40&320&14&52&48{,}06&\cdot\ \cdot\ \cdot\ \cdot&53568{,}06\\
50&310&18&12&49{,}21&\cdot\ \cdot\ \cdot\ \cdot&65569{,}21\\
60&300&21&16&17{,}02&\cdot\ \cdot\ \cdot\ \cdot&76577{,}02\\
70&290&23&58&42{,}92&\cdot\ \cdot\ \cdot\ \cdot&86322{,}92\\
80&280&26&14&55{,}85&\cdot\ \cdot\ \cdot\ \cdot&94495{,}85\\
90&270&27&58&52{,}36&\cdot\ \cdot\ \cdot\ \cdot&100732{,}36\\
100&260&29&3&28{,}13&\cdot\ \cdot\ \cdot\ \cdot&104608{,}13\\
110&250&29&20&33{,}68&\cdot\ \cdot\ \cdot\ \cdot&105633{,}68\\
120&240&28&41&2{,}10&\cdot\ \cdot\ \cdot\ \cdot&103262{,}10\\
130&230&26&55&22{,}77&\cdot\ \cdot\ \cdot\ \cdot&96922{,}77\\
140&220&23&55&2{,}70&\cdot\ \cdot\ \cdot\ \cdot&86102{,}70\\
150&210&19&35&0{,}79&\cdot\ \cdot\ \cdot\ \cdot&70500{,}79\\
160&200&13&57&40{,}52&\cdot\ \cdot\ \cdot\ \cdot&50260{,}52\\
170&190&7&16&58{,}33&\cdot\ \cdot\ \cdot\ \cdot&26218{,}33\\
180&180&0&&&&0
\end{array}
\]

\Needspace{15\baselineskip}
We divide this period of \(36\) terms into six smaller ones of six terms each; the six values of the function \(X\) contained in each period will be represented by a formula of the form
\[
\begin{aligned}
&\gamma+\gamma'\cos x+\gamma''\cos 2x+\gamma'''\cos 3x\\[6pt]
&\quad+\delta'\sin x+\delta''\sin 2x+\delta'''\sin 3x
\end{aligned}
\]
\SourceLeaf{335}
where we find the following values for the coefficients \(\gamma,\gamma',\delta'\) etc.:


\begin{landscape}\thispagestyle{empty}\begingroup\fontsize{9}{12}\selectfont\setlength{\arraycolsep}{3pt}\renewcommand{\arraystretch}{1.3}
\[
\begin{array}{l|r|r|r|r|r|r|r|r}
\text{Period}&y=6x&\gamma&\gamma'&\gamma''&\gamma'''&\delta'&\delta''&\delta'''\\ \hline
\text{first}&0^{\mathord{0}}&0&0&0&0&
-103830{,}165&
+15406{,}638&0\\
\text{second}&60&+56''205&-217''757&+761''671&-1528{,}825&-103937{,}346&+15841{,}387&-882{,}667\\
\text{third}&120&+56{,}115&-217{,}467&+760{,}805&-1527{,}397&-104151{,}314&+16709{,}402&-2645{,}530\\
\text{fourth}&180&0&0&0&0&-104258{,}100&+17142{,}684&-3525{,}740\\
\text{fifth}&240&-56{,}115&-217{,}467&-760{,}805&+1527{,}397&-104151{,}314&+16709{,}402&-2645{,}530\\
\text{sixth}&300&-56{,}205&-217{,}757&-761{,}671&+1528{,}825&-103937{,}346&+15841{,}387&-882{,}667
\end{array}
\]
\endgroup\end{landscape}


\Needspace{39\baselineskip}
The individual coefficients \(\gamma,\gamma',\gamma'',\gamma''',\delta',\delta'',\delta'''\) will in turn be reduced to a form such as
\[
\begin{aligned}
&\varepsilon+\varepsilon'\cos 6x+\varepsilon''\cos 12x+\varepsilon'''\cos 18x\\[6pt]
&\quad+\zeta'\sin 6x+\zeta''\sin 12x+\zeta'''\sin 18x
\end{aligned}
\]
but it is readily seen that, for the first four, \(\varepsilon,\varepsilon',\varepsilon'',\varepsilon'''\) must vanish; and for the last three, \(\zeta',\zeta'',\zeta'''\). In this way one obtains
\[
\begin{aligned}
&\gamma=+64''{,}848\sin 6x+0''{,}052\sin 12x\\[6pt]
&\gamma'=-251''{,}277\sin 6x-0''{,}167\sin 12x\\[6pt]
&\gamma''=+879''{,}002\sin 6x+0''{,}500\sin 12x\\[6pt]
&\gamma'''=-1764''{,}511\sin 6x-0''{,}824\sin 12x\\[6pt]
&\delta'=-104044''{,}264+213''{,}968\cos 6x+0''{,}132\cos 12x\\[6pt]
&\qquad{}+0''{,}000\cos 18x\\[6pt]
&\delta''=+16275''{,}150-868''{,}020\cos 6x-0''{,}489\cos 12x\\[6pt]
&\qquad{}-0''{,}003\cos 18x\\[6pt]
&\delta'''=-1763''{,}689+1762''{,}868\cos 6x+0''{,}819\cos 12x\\[6pt]
&\qquad{}+0''{,}002\cos 18x
\end{aligned}
\]

\Needspace{24\baselineskip}
On substituting these values for \(\gamma,\gamma'\) etc. and observing the rules of the preceding article, the following function results for the equation of the center, in which the individual coefficients are accurate to within one hundredth of a second of arc.


\[
\begin{array}{r|r}
-104044''264\sin x&-1''643\sin 9x\\
+16275{,}150\sin 2x&+0{,}494\sin 10x\\
-3527{,}378\sin 3x&-0{,}149\sin 11x\\
+873{,}511\sin 4x&+0{,}052\sin 12x\\
-232{,}622\sin 5x&-0{,}017\sin 13x\\
+64{,}848\sin 6x&+0{,}006\sin 14x\\
-18{,}655\sin 7x&-0{,}004\sin 15x\\
+5{,}491\sin 8x&+0{,}003\sin 16x
\end{array}
\]


\Needspace{36\baselineskip}
\SourceLeaf{336}
In this way the calculation evidently proceeds as if, in place of the arc \(x\), another, \(x'=x-a\), were introduced. Often, too, the calculation can be greatly eased if, in place of the function \(X\), another, \(X-\mathfrak{x}\), is used, where for \(\mathfrak{x}\) any function similar to \(X\) may be taken, provided it does not exceed the \(\frac12\mu-\frac12{}^{\mathrm{th}}\) order for an odd value or the \(\frac12\mu^{\mathrm{th}}\) order for an even value of \(\mu\), and in the latter case the coefficients of the terms \(\cos\frac12\mu x\) and \(\sin\frac12\mu x\) (if it rises that far) are in the required ratio \((\cos\frac12\mu a:\sin\frac12\mu a)\). It is evident that, if the values of the function \(\mathfrak{x}\) for \(x=a.b,c,d\) etc. are respectively \(\mathfrak{A},\mathfrak{B},\mathfrak{C},\mathfrak{D}\) etc., the values of the function \(X-\mathfrak{x}\) for the same values will be respectively \(A-\mathfrak{A},B-\mathfrak{B},C-\mathfrak{C},D-\mathfrak{D}\) etc.; if, further, the function \(Z\) satisfying these values is found by the preceding rules, the function \(\mathfrak{x}+Z\) will be identical with that which would result from applying these rules directly to the values of the function \(X\) itself. The function \(\mathfrak{x}\) must always be determined so that the proposed values of the function \(X\) are closely represented by it, whereby the quantities \(A-\mathfrak{A},B-\mathfrak{B},C-\mathfrak{C},D-\mathfrak{D}\) etc. will become small and more manageable in the calculation. Moreover, for \(\mathfrak{x}\) a constant quantity may also be taken, e.g. this one, \(\frac1\mu(A+B+C+D+\text{etc.})\), which constitutes the invariable part of the function sought.

\begin{center}\rule{0.28\linewidth}{0.4pt}\end{center}
\endgroup
